\documentclass[a3paper]{article}
\usepackage[margin=0.5cm]{geometry}
\usepackage{graphicx}
\usepackage{tikz}
\usepackage{xifthen}

\usepackage{etoolbox}
\usepackage{pgfmath} % para \pgfmathtruncatemacro
\usepackage[T1]{fontenc}
% \usepackage[OT1]{fontenc}
\usetikzlibrary{calc}
\usetikzlibrary{positioning}
\usepackage{pgffor}
\pagestyle{empty}
        
% === CONFIGURACIÓN DE CARTAS ===
\newlength{\cardwidth}\setlength{\cardwidth}{5.7cm}
\newlength{\cardheight}\setlength{\cardheight}{9cm}

% Macro para dibujar una carta
% \newcommand{\drawcard}[4]{%
%   % #1 = tipo (macro/micro), #2 = nombre, #3 = x, #4 = y
%   \begin{scope}[shift={(#3,#4)}]
%     \draw[rounded corners=5pt,thick] (0,0) rectangle (\cardwidth,-\cardheight);
%     \node[anchor=center] at (\cardwidth/2,-\cardheight/2)
%       {\includegraphics[width=0.9\cardwidth,height=0.9\cardheight,keepaspectratio]{#1/#2.png}};
%     \ifx#1micro
%       \node[anchor=south east,xshift=-2mm,yshift=2mm] at (\cardwidth,-\cardheight)
%         {\includegraphics[width=1.5cm]{macro/#2.png}};
%     \fi
%   \end{scope}
% }

\newcommand{\drawmicro}[3]{%
  % #1 = nombre archivo
  % #2 = nombre quimico
  % #3 = x, #4 = y (coordenadas en cm)
  \begin{scope}[shift={#3}]
    % Marco de la carta
    \draw[rounded corners=5pt,thick] (0,0) rectangle (\cardwidth,-\cardheight);

    \node[anchor=center, shift={(0,1)}] at (.5\cardwidth,-.2\cardheight)
      {\LARGE #2};

    % Imagen principal
    \node[anchor=center, shift={(0,.2)}] at (\cardwidth/2,-\cardheight/2)
      {\includegraphics[width=0.9\cardwidth,height=0.9\cardheight,keepaspectratio]{micro/#1.png}};

    % Definir centro y radio del círculo
    \begin{scope}
      % Posicionar el centro en esquina inferior derecha
      \clip (\cardwidth-1.5cm,-\cardheight+1.5cm) circle[radius=1.4cm];
      % Insertar la macro centrada
      \node[anchor=center] at (\cardwidth-1.5cm,-\cardheight+1.5cm)
         {\includegraphics[width=0.9\cardwidth,height=0.9\cardheight,keepaspectratio]{inset/#1.png}};
    \end{scope}
    % Opcional: dibujar borde del círculo
    \draw[thick] (\cardwidth-1.5cm,-\cardheight+1.5cm) circle[radius=1.4cm];
  \end{scope}
}
\newcommand{\drawmacro}[2]{%
  % #1 = nombre
  % #2 = x, #3 = y (coordenadas en cm)
  \begin{scope}[shift={#2}]
    % Marco de la carta
    \draw[rounded corners=5pt,thick] (0,0) rectangle (\cardwidth,-\cardheight);

    % Imagen principal
    \node[anchor=center] at (\cardwidth/2,-\cardheight/2)
      {\includegraphics[width=0.9\cardwidth,height=0.9\cardheight,keepaspectratio]{macro/#1.png}};

  \end{scope}
}


\begin{document}

\gdef\globaln{0}
\gdef\globalcw{5.7} 
\gdef\globalch{9} 

% \def\molecules{{cafeina,agua,etanol,glucosa,urea,acetona,metanol}}
% El origen
\begin{tikzpicture}[remember picture, overlay]
    \coordinate (origen) at (0,0);
\end{tikzpicture}

\foreach \mol/\name in {
ac-acetilsalicilico/ácido acetilsalicílico
,ac-ascorbico/ácido ascórbico
,adn/ADN
,agua/agua
,alicina/alicina
,cafeina/cafeína
,celulosa/celulosa
,clorofila-a/clorofila A
,etanol/etanol
,ftalocianina-Cu/ftalocianina de cobre
,glucosa/glucosa
,hidroxiapatita/hidroxiapatita
,laurilsulfato-Na/laurilsulfato de sodio
,lignina/lignina
,metano/metano
,nicotina/nicotina
,queratina/queratina
,SiO2/dióxido de silicio
,sucrosa/sucrosa
,sulfanona/sulfanona
   ,ac-acetilsalicilico/ácido acetilsalicílico
,ac-ascorbico/ácido ascórbico
,adn/ADN
,agua/agua
,alicina/alicina
,cafeina/cafeína
,celulosa/celulosa
,clorofila-a/clorofila A
,etanol/etanol
,ftalocianina-Cu/ftalocianina de cobre
,glucosa/glucosa
,hidroxiapatita/hidroxiapatita
,laurilsulfato-Na/laurilsulfato de sodio
,lignina/lignina
,metano/metano
,nicotina/nicotina
,queratina/queratina
,SiO2/dióxido de silicio
,sucrosa/sucrosa
,sulfanona/sulfanona 
}{

    % Calcula fila y columna de la pareja (macro+micro)
    \pgfmathtruncatemacro{\col}{int(mod(\globaln,2))}
    \pgfmathtruncatemacro{\row}{int(\globaln/2)}
    
    % Coordenadas
    \pgfmathsetmacro{\x}{\col*2*(\globalcw+0.5)}
    \pgfmathsetmacro{\y}{-\row*(\globalch+0.5)}
    
    \begin{tikzpicture}[remember picture,overlay]
        \pgfmathtruncatemacro\xc{\globaln+1}
        \path(origen) ++(\x,\y) coordinate(A);
        \path(A) ++(\globalcw,0) ++(0.5,0) coordinate(B);
        \drawmacro{\mol}{(A)}
        \drawmicro{\mol}{\name}{(B)}
        \global\let\globaln=\xc
    \end{tikzpicture}
    
    % \advance\c by 1
    % % Salto de página si se llenó la hoja
    \ifnum\globaln=10
      \clearpage
      \begin{tikzpicture}[remember picture, overlay]
          \coordinate (origen) at (0,0);
      \end{tikzpicture}
      \pgfmathtruncatemacro\xc{0}
      \global\let\globaln=\xc
    \fi
}

\end{document}

